\chapter{刻度的测量}

记号约定:
\begin{itemize}
	\item $\left(E,\leq\right)$是全序集,它有最小元$\omega$.
	\item $I$是$E$中的下集,并且包含$\omega$.令$I_{+}=\left\{x\in I\middle|x>\omega\right\}$.
	\item $E$上部分定义了局部定义的乘法$\ast:I\times I\to E$.
\end{itemize}

\begin{definition}{$I$中的乘积}{}
	设$\left(x_{i}\right)_{i=1}^{p}$是$I$中的有限序列,递归地定义$\prod\limits_{i=1}^{p}x_{i}$如下:
	\begin{enumerate}
		\item 当$p=0$时$\prod\limits_{i=1}^{p}x_{i}=\omega$,当$p=1$时$\prod\limits_{i=1}^{p}x_{i}=x_{1}$,当$p=2$时$\prod\limits_{i=1}^{p}x_{i}=x_{1}\ast x_{2}$.
		\item 当$p\geq 3$,$\prod\limits_{i=1}^{p-1}x_{i}$有定义且属于$I$时,$\prod\limits_{i=1}^{p}x_{i}=\left(\prod\limits_{i=1}^{p-1}x_{i}\right)\ast x_{p}$.
	\end{enumerate}
\end{definition}

\begin{definition}{$I$中元素的幂}{}
	设$x\in I$.对每个$p\in\N$记$x^{p}\coloneq\prod\limits_{i=1}^{p}x$.
\end{definition}

考虑如下条件:
\begin{description}
	\item ($\mathrm{GR}_{1}$)\textbf{$\omega$是局部单位元且$\ast$满足局部结合律:}如下条件同时成立:
	\begin{itemize}
		\item $\forall x\in I\left(xw=wx=x\right)$.
		\item $\forall x,y,z\in I\left(x\ast y,y\ast z,\left(x\ast y\right)\ast z,x\ast\left(y\ast z\right)\in I\implies\left(x\ast y\right)\ast z=x\ast\left(y\ast z\right)\right)$.
	\end{itemize}
	\item ($\mathrm{GR}_{2}$)\textbf{$\ast$和$\leq$相容:}$\forall x,y\in I\left(x<y\implies \forall z\in I\left(xz<yz\wedge zx<zy\right)\right)$.
	\item ($\mathrm{GR}_{3}$)\textbf{$I$无最小正量且``可任意细分'':}$I_{+}$非空且无最小元;$\forall x,y\in I\left(x<y\implies\exists z\in I_{w}\left(xz<y\right)\right)$.
	\item ($\mathrm{GR}_{3a}$)\textbf{$I$无最小正量且``可精确作差'':}$I_{+}$非空且无最小元;$\forall x,y\in I\left(x<y\implies\exists z\in I\left(xz=y\right)\right)$.
	\item ($\mathrm{GR}_{4}$)\textbf{Archimedes公理:}$\forall x\in I_{+}\forall y\in I\exists n\geq 1\left(x^{n}\in I\wedge x^{n}>y\right)$.
	\item ($\mathrm{GR}_{4a}$)\textbf{序列完备性:}$I$中每个单调递增且有上界的序列$\left(x_{n}\right)_{n=0}^{\infty}$在$I$中有上确界.
\end{description}

\begin{proposition}{$I$中的不等式可以逐项相差}{}
	设($\mathrm{GR}_{2}$)成立.
	\begin{enumerate}
		\item $\forall x,x^{\prime},y,y^{\prime}\in I\left(x<y\wedge x^{\prime}<y^{\prime}\wedge x\ast x^{\prime}\in I\wedge y\ast y^{\prime}\in I\implies x\ast x^{\prime}<y\ast y^{\prime}\right)$.
		\item 如果($\mathrm{GR}_{1}$)成立,那么$\forall x\in I_{+}\forall y\in I\left(y<yx\right)$.
	\end{enumerate}
\end{proposition}

\begin{proposition}{}{}
	设$x,y\in I$.
	\begin{enumerate}
		\item 设$p\geq 3$.如果$x^{p}$有定义,那么对每个$2\leq q\leq p$,$x^{q}$有定义且属于$I$.
		\item 设$p\geq 3$,($\mathrm{GR}_{1}$)和($\mathrm{GR}_{2}$)成立.如果$x\in I_{+}$且$x^{p}$有定义,那么对每个$1\leq q\leq p-1$有$\omega<x^{q}<x^{p}$.
		\item 设($\mathrm{GR}_{2}$)成立,$p\geq 3$.如果$x<y$且$y^{p}$有定义,那么$x^{p}$有定义且$x^{p}<y^{p}$.
		\item 设($\mathrm{GR}_{1}$)成立.
		\begin{itemize}
			\item 如果$m,n\in\N$且$x^{m+n}$有定义,那么$x^{m}\ast x^{n}$有定义且$x^{m+n}=x^{m}\ast x^{n}$.
			\item 设($\mathrm{GR}_{2}$)成立.如果$m,n\in\N$,$x^{m}\ast x^{n}$有定义且属于$I$,那么$x^{m+n}$有定义且$x^{m+n}=x^{m}\ast x^{n}$.
		\end{itemize}
		\item 设($\mathrm{GR}_{1}$)成立.
		\begin{itemize}
			\item 如果$m,n\in\N$且$x^{mn}$有定义,那么$\left(x^{m}\right)^{n}$有定义且$x^{mn}=\left(x^{m}\right)^{n}$.
			\item 设($\mathrm{GR}_{2}$)成立.如果$m,n\in\N$,$\left(x^{m}\right)^{n}$有定义且属于$I$,那么$x^{mn}$有定义且$x^{mn}=\left(x^{m}\right)^{n}$.
		\end{itemize}
	\end{enumerate}
\end{proposition}

\begin{proposition}{}{}
	设($\mathrm{GR}_{2}$),($\mathrm{GR}_{3}$)成立.
	\begin{enumerate}
		\item $\forall x\in I_{+}\exists y\in I_{+}\left(y^{2}\leq x\right)$.
		\item 如果($\mathrm{GR}_{1}$)成立,那么$\forall n\in\N\forall x\in I_{+}\exists u\in I_{+}\left(u^{2^{n}}\leq x\right)$
	\end{enumerate}
\end{proposition}

\begin{definition}{测量系}{}
	当($\mathrm{GR}_{1}$),($\mathrm{GR}_{2}$),($\mathrm{GR}_{3}$)成立时,称$\left(E,\leq,\omega,I,\ast\right)$是测量系.
\end{definition}

\begin{definition}{近似商}{}
	设$x\in I,y\in I\setminus\left\{\omega\right\}$.定义$\left(x:y\right)\coloneq\max\left\{n\in\N\middle|y^{n}\text{有定义},y^{n}\leq x\right\}$.
\end{definition}

\begin{proposition}{近似上的次可加性和链式不等式}{}
	设($\mathrm{GR}_{1}$),($\mathrm{GR}_{2}$),($\mathrm{GR}_{4}$)成立.如果$x,y,z\in I,xy\in I,z\neq\omega$,那么
	\begin{gather*}
		\left(x:z\right)+\left(y:z\right)\leq\left(x\ast y:z\right)\leq\left(x:z\right)+\left(y:z\right)+1,\\
		\left(x:y\right)\left(y:z\right)\leq\left(x:z\right);\\
		\left(\left(x:y\right)+1\right)\left(\left(y:z\right)+1\right)\geq\left(x:z\right)+1.
	\end{gather*}
\end{proposition}

\begin{proposition}{}{}
	设$\left(E,\leq,\omega,I,\ast\right)$是测量系,($\mathrm{GR}_{4}$)成立.令$\mathfrak{F}$是$I_{+}$按$\geq$所成定向集的截面滤子.
	\begin{enumerate}
		\item 对每个$u\in I_{+}$和$z\in I\setminus\left\{\omega\right\}$,$\left(u:z\right)$沿$\mathfrak{F}$收敛到$+\infty$.
		\item 对每个$a\in I_{+}$和$x\in I$,映射$f_{a,x}:I\cap\left(-\infty,a\right]\to\R,z\mapsto\frac{\left(x:z\right)}{\left(a:z\right)}$满足如下条件:
		\begin{itemize}
			\item $f_{a,x}$沿$\mathfrak{F}$收敛.
			\item 当$x=\omega$时,$f$收敛到$0$.
			\item 当$x\in I_{+}$时,$\im\left(f\right)$是$\R_{>0}^{\times}$的Cauchy滤子基,并且收敛到一个正实数.
		\end{itemize}
	\end{enumerate}
\end{proposition}

\begin{proposition}{测度函数的性质}{}
	设$\left(E,\leq,\omega,I,\ast\right)$是测量系,($\mathrm{GR}_{4}$)成立,$\mathfrak{F}$是$I_{+}$按$\geq$所成定向集的截面滤子,$a\in I_{+}$.令
	\begin{align*}
		f_{a}:I\to\R,x\mapsto\limit_{\mathfrak{F}}\frac{\left(x:z\right)}{\left(a:z\right)}.
	\end{align*}
	\begin{enumerate}
		\item $f\left(\omega\right)=0,f\left(I_{+}\right)\subseteq\R_{>0},f_{a}\left(a\right)=1$.
		\item $\forall x,y\in I\left(x\ast y\in I\implies f\left(x\ast y\right)=f\left(x\right)+f\left(y\right)\right)$.
		\item $\forall x,y\in I\left(x<y\implies f\left(x\right)<f\left(y\right)\right)$.
		\item 对每个$b\in I$,集合$\im\left(f\right)\cap\left[0,f\left(b\right)\right]$在$\left[0,f\left(b\right)\right]$中稠密.
	\end{enumerate}
\end{proposition}

\begin{remark}{}{}
	\begin{enumerate}
		\item 在这一命题的条件下,$\forall x,y\in I\left(x\ast y=y\ast x\right)$.
		\item 如果$g:I\to\R$满足如下条件:
		\begin{itemize}
			\item $\forall x,y\in I\left(x\ast y\in I\implies g\left(x\ast y\right)=g\left(x\right)+g\left(y\right)\right)$,
			\item $\forall x,y\in I\left(x<y\implies g\left(x\right)<g\left(y\right)\right)$,
		\end{itemize}
		那么存在$\lambda\in\R_{>0}$使得$g=\lambda f$.
	\end{enumerate}
\end{remark}

\begin{proposition}{}{}
	如果($\mathrm{GR}_{1}$),($\mathrm{GR}_{2}$),($\mathrm{GR}_{3a}$),($\mathrm{GR}_{4a}$)成立,那么($\mathrm{GR}_{4}$)成立.
\end{proposition}

\begin{proposition}{}{}
	如果($\mathrm{GR}_{1}$),($\mathrm{GR}_{2}$),($\mathrm{GR}_{3a}$),($\mathrm{GR}_{4a}$)成立,那么存在$\alpha\in\left(0,+\infty\right]$和严格递增映射$f:I\to\left[0,\alpha\right)$满足如下条件:
	\begin{enumerate}
		\item $f\left(\omega\right)=0$;
		\item $\forall x,y\in I\left(x\ast y\in I\implies f\left(x\ast y\right)=f\left(x\right)+f\left(y\right)\right)$;
		\item $f$是满射.
	\end{enumerate}
	特别地,如果$I=E$,那么$\alpha=+\infty$.
\end{proposition}